\documentclass[10pt,a4paper]{amsart}
\usepackage[margin=19mm]{geometry}
\usepackage{amsmath,amssymb,mathtools}
\usepackage[scr=rsfs,cal=euler]{mathalfa}
\usepackage{xfrac}
\usepackage[unicode,hidelinks,bookmarks=false]{hyperref}
\newcommand{\ie}{i.e.\ }
\newcommand{\cf}{cf.\ }
\newcommand{\resp}{respectively\ }
\newcommand{\Subsection}{\subsection}
\providecommand{\nobreakdash}{\nobreak\hbox{-}}
\setlength{\emergencystretch}{2em}
\multlinegap0pt
\usepackage{amssymb}
\usepackage{bm}
\usepackage{float}
\usepackage{xcolor}
\usepackage{algorithm}\usepackage{algpseudocode}
\newtheorem{lemma}{Lemma}
\newtheorem{theorem}{Theorem}
\title{Stable High-Order Interpolation on the Grassmann Manifold by Maximum-Volume Coordinates and Arnoldi Orthogonalization}
\author{Qiang Niu, Wen Jiang, Jie Fei, Ruoyu Xiong, Yuxuan Li}
\date{Source: arXiv:2603.23858v2 (CC BY 4.0)}
\begin{document}
\maketitle

\noindent\emph{Source and licence.} Stable High-Order Interpolation on the Grassmann Manifold by Maximum-Volume Coordinates and Arnoldi Orthogonalization. Version arXiv:2603.23858v2 (CC BY 4.0), distributed by arXiv under the Creative Commons Attribution 4.0 International licence, http://creativecommons.org/licenses/by/4.0/, as recorded in the arXiv licence metadata for this exact version and shown on the abstract page https://arxiv.org/abs/2603.23858v2. Full licence notice: \texttt{licenses/arxiv-2603.23858-CC-BY-4.0.txt}.\par\smallskip

\noindent\textit{Editorial scope.} Counted unit: the article's theorem thm:err\_prop (source0.tex lines 549--555). The article's complete original proof of it is delivered with it (source0.tex lines 557--617, one \texttt{proof} environment of 61 lines). No mathematics is written, repaired or re-derived: the statement and the proof are the article's own lines, in the article's own notation, and no retained line is substituted or re-flowed.  Retained uncounted as the source-local context the proof needs: lines 519--521.  Nothing was omitted from any retained span, so the package carries no omission record.  No retained line contains a citation, so the wrapper carries no bibliography and no bibliography entry.  No retained line contains a figure environment, an image-inclusion command or drawing code, so no image is delivered and the package declares no asset dependency.  Editorial formatting only, in addition: an AMS \texttt{amsart} preamble that restates the article's own declarations and macros used by the retained lines, namely the environments \texttt{lemma}, \texttt{theorem} on separate counters, together with the standard packages those lines need.  The retained environments print in this document as Lemma 1, Theorem 1 in order of appearance, whereas the article prints the same items with the numbering of its own full document.  The complete native source of the article as distributed in the arXiv e-print for this exact version is bundled unchanged as \texttt{sources/197114/source0.tex}.
\begin{lemma} \label{lem:smooth_iso}
The mapping $\Xi \mapsto \widetilde{U}(\Xi)$ defined in Lemma 1 is $C^\infty$ on its domain. Furthermore, at $\Xi = 0$, the differential $d\widetilde{U}|_0(\Delta) = \begin{bmatrix} 0 \\ \Delta \end{bmatrix}$, which constitutes an isometry between the tangent space $T_0 \mathbb{R}^{(n-p) \times p}$ and the tangent space $T_{\widetilde{U}(0)} \text{St}(n,p)$, where the base point is $\widetilde{U}(0) = \begin{bmatrix} I_p \\ 0 \end{bmatrix}$.
\end{lemma}

\begin{theorem} \label{thm:err_prop}
Let $\Xi$ be the true MV coordinates and $\hat{\Xi} = \Xi + \Delta$ be the interpolated coordinates. Under the Householder-stabilized system, the reconstruction error for the orthogonal projector $P = UU^T$ satisfies
\begin{equation} \label{eq:err_bound_final}
\|\hat{P} - P\|_F \leq \sqrt{2} \cdot \|\widetilde U_1\|_2 \cdot \|\Delta\|_F + \mathcal{O}(\|\Xi\|_2 \|\Delta\|_F + \|\Delta\|_F^2),
\end{equation}
where $\|\widetilde U_1\|_2$ is the spectral norm of the upper $p \times p$ block of the reconstructed Stiefel representative.
\end{theorem}

\begin{proof}
Let $\Psi(\Xi) = \widetilde{U}(\Xi)\widetilde{U}(\Xi)^T$ be the mapping from the coordinate domain to the Grassmannian. Define the line segment $\gamma(t) = \Xi + t\Delta$ for $t \in [0,1]$. By the Mean Value Theorem for matrix-valued functions, we have
 $$
\hat{P} - P = \Psi(\Xi + \Delta) - \Psi(\Xi) = \int_0^1 d\Psi(\gamma(t))[\Delta] \, dt.
 $$
Taking the Frobenius norm on both sides yields
\begin{equation}
\|\hat{P} - P\|_F \leq \int_0^1 \left\| d\Psi(\gamma(t))[\Delta] \right\|_F dt \leq \sup_{t \in [0,1]} \left\| d\Psi(\gamma(t))[\Delta] \right\|_F. \label{eq:norm_integral}
\end{equation}
To bound the differential norm at any point $X \in \{\gamma(t)\}$, we evaluate $d\Psi(X)[\Delta] = (d\widetilde U)\widetilde U^T + \widetilde U(d\widetilde U)^T$. Using the reconstruction formula $\widetilde U(X) = \begin{bmatrix} I_p \\ X \end{bmatrix} S(X)$ where $S(X) = (I_p + X^T X)^{-1/2}$, the differential is
\begin{equation} \label{eq:du_exact_rigorous}
d\widetilde U = D_1 + D_2.
\end{equation}
where  $D_1 = \begin{bmatrix} 0 \\ \Delta \end{bmatrix} S$ and $D_2 = \begin{bmatrix} I_p \\ X \end{bmatrix} dS$.
Differentiating both sides of $S^{-2} = I_p + X^T X$ yields:
$$ d(S^{-2}) = \Delta^T X + X^T \Delta. $$
Applying the product rule to $S^{-2} = S^{-1}S^{-1}$ and using the matrix differential identity $d(S^{-1}) = -S^{-1}(dS)S^{-1}$, we expand the left side as
$ d(S^{-2})  = -S^{-1}(dS)S^{-2} - S^{-2}(dS)S^{-1}. $
Equating the two expressions and pre- and post-multiplying by $S$ give the Sylvester equation for $dS$:
$$ (dS)S^{-1} + S^{-1}(dS) = -S(\Delta^T X + X^T \Delta)S. $$
Its unique solution is given by the integral representation
$$ dS = \int_0^\infty e^{-S^{-1}\tau} [S(\Delta^T X + X^T \Delta)S] e^{-S^{-1}\tau} d\tau. $$
Taking the spectral norm and applying the bound $|e^{-S^{-1}\tau}|_2 \le e^{-\tau/|S|_2}$, the scalar integral can be evaluated to produce
$$ \|dS\|_2 \le \frac{\|S\|_2}{2} \|S(\Delta^T X + X^T \Delta)S\|_2 \le \|S\|_2^3 \|X\|_2 \|\Delta\|_F. $$
Substituting $d\tilde{U}$ into $d\Psi$, we decompose it into a principal part and a remainder
$$ d\Psi = (D_1 \tilde{U}^T + \tilde{U} D_1^T) + \underbrace{(D_2 \tilde{U}^T + \tilde{U} D_2^T)}_{\mathcal{O}(\|X\|_2\|\Delta\|_F)}. $$
Let $A = D_1 \tilde{U}^T = \begin{bmatrix} 0 \\ \Delta S \end{bmatrix} \tilde{U}^T$. Applying the triangle inequality to the remainder $\mathcal{R} = D_2 \tilde{U}^T + \tilde{U} D_2^T$ gives 
$$ \|\mathcal{R}\|_F \le 2 \|D_2\|_F \|\tilde{U}\|_2 \le 2 \left|\left| \begin{bmatrix} I_p \\ X \end{bmatrix} \right|\right|_2 \|dS\|_F \le 2\sqrt{p} \sqrt{1 + \|X\|_2^2} \cdot \|S\|_2^3 \|X\|_2 \|\Delta\|_F. $$
Thus the remainder is rigorously bounded by $\|\mathcal{R}\|_F = \mathcal{O}(\|X\|_2\|\Delta\|_F)$.  Next, we evaluate $\|A + A^T\|_F$. By the properties of the Frobenius norm,
\begin{equation}\label{eq:a_plus_at_norm}
\|A + A^T\|_F^2 = \|A\|_F^2 + \|A^T\|_F^2 + 2\operatorname{tr}(A^2) = 2\|A\|_F^2 + 2\operatorname{tr}(A^2).
\end{equation}
Since $\widetilde U$ has orthonormal columns ($\widetilde U^T \widetilde U = I_p$), we have $$\|A\|_F^2 = \operatorname{tr}(A^T A) = \operatorname{tr}(\widetilde U D_1^T D_1 \widetilde U^T) = \operatorname{tr}(D_1^T D_1) = \|D_1\|_F^2 = \|\Delta S\|_F^2.$$
Moreover, 
 $$
A^2 = \left( \begin{bmatrix} 0 \\ \Delta S \end{bmatrix} \widetilde U^T \right)^2 = \begin{bmatrix} 0 \\ \Delta S \end{bmatrix} \left( \begin{bmatrix} S & S X^T \end{bmatrix} \begin{bmatrix} 0 \\ \Delta S \end{bmatrix} \right) \widetilde U^T = \begin{bmatrix} 0 \\ \Delta S^2 X^T \Delta S \end{bmatrix} \widetilde U^T.
 $$
Thus, $\operatorname{tr}(A^2) = \operatorname{tr}(S X^T \Delta S^2 X^T \Delta S)$. By Cauchy-Schwarz inequality for trace inner products $|\operatorname{tr}(Y Z)| \leq \|Y\|_F \|Z\|_F$, letting $Y = S X^T \Delta S$ and $Z = S X^T \Delta S$, we get
 $$
|\operatorname{tr}(A^2)| \leq \|S X^T \Delta S\|_F^2 \leq \|S\|_2^4 \|X\|_2^2 \|\Delta\|_F^2.
 $$
Substituting back into \eqref{eq:a_plus_at_norm}, we obtain
 $$
\|A + A^T\|_F = \sqrt{2\|\Delta S\|_F^2 + \mathcal{O}(\|X\|_2^2 \|\Delta\|_F^2)} = \sqrt{2} \|\Delta S\|_F + \mathcal{O}(\|X\|_2  \|\Delta\|_F).
 $$
Using the sub-multiplicative property $\|\Delta S\|_F \leq \|S\|_2 \|\Delta\|_F$,
 $$
\|A + A^T\|_F \leq \sqrt{2} \|S\|_2 \|\Delta\|_F + \mathcal{O}(\|X\|_2  \|\Delta\|_F).
 $$
Combining the bounds for the principal part and the remainder $\mathcal{R}$, the differential norm at $X$ is bounded by
 $$
\|d\Psi(X)[\Delta]\|_F \leq \sqrt{2} \|S(X)\|_2 \|\Delta\|_F + \mathcal{O}(\|X\|_2   \|\Delta\|_F).
 $$
Returning to the supremum in \eqref{eq:norm_integral}, since $X \in \{\Xi + t\Delta\}$, we have $\|X\|_2 \leq \|\Xi\|_2 + \|\Delta\|_2$.
Furthermore, as established in Lemma~\ref{lem:smooth_iso}, the matrix function $S(\cdot)$ is smooth, ensuring $\|S(X)\|_2 = \|S(\Xi)\|_2 + \mathcal{O}(\|\Delta\|_F)$.
Since the upper block of $\widetilde U(\Xi)$ is $S(\Xi)$, we have $\|\widetilde U_1\|_2 = \|S(\Xi)\|_2$. Substituting these limits into the integral bound completes the strict derivation
$$
\|\hat{P} - P\|_F \leq \sqrt{2} \|\widetilde U_1\|_2 \|\Delta\|_F + \mathcal{O}(\|\Xi\|_2 \|\Delta\|_F + \|\Delta\|_F^2).
$$
The constant $\sqrt{2}$ arises from the geometry of the projector space, while the spectral norm $\|\widetilde U_1\|_2$ ensures that the reconstruction mapping does not amplify the coordinate-domain error.
\end{proof}

\end{document}
